\section{Results}
\label{sec:results}

\subsection{Main Results: Binary Entropy by Strategy and Dataset}

\tabref{tab:entropy_gpt41} and \tabref{tab:entropy_mini} present mean binary entropy across all dataset--strategy combinations for \gptfull and \gptmini, respectively.

\begin{table}[t]
    \centering
    \caption{Binary entropy (mean $\pm$ std) for \gptfull across prompt strategies. Higher is better; 1.000 indicates a perfect half-split. Best result per dataset in \textbf{bold}.}
    \resizebox{\textwidth}{!}{%
    \begin{tabular}{@{}lccc@{}}
        \toprule
        \textbf{Dataset} & \textbf{\basic} & \textbf{\explicit} & \textbf{\chainofthought} \\
        \midrule
        Mixed Categories & 0.841 {\scriptsize $\pm$ 0.365} & \textbf{1.000} {\scriptsize $\pm$ 0.000} & \textbf{1.000} {\scriptsize $\pm$ 0.000} \\
        \mcrae 8 & 0.206 {\scriptsize $\pm$ 0.367} & 0.997 {\scriptsize $\pm$ 0.012} & \textbf{1.000} {\scriptsize $\pm$ 0.000} \\
        GPT 8 & 0.177 {\scriptsize $\pm$ 0.375} & \textbf{0.998} {\scriptsize $\pm$ 0.010} & \textbf{0.998} {\scriptsize $\pm$ 0.010} \\
        \mcrae 16 & 0.484 {\scriptsize $\pm$ 0.475} & 0.995 {\scriptsize $\pm$ 0.014} & \textbf{1.000} {\scriptsize $\pm$ 0.000} \\
        \bigbench & \textbf{0.992} {\scriptsize $\pm$ 0.000} & 0.998 {\scriptsize $\pm$ 0.002} & 0.942 {\scriptsize $\pm$ 0.107} \\
        \things 8 & 0.410 {\scriptsize $\pm$ 0.394} & 0.932 {\scriptsize $\pm$ 0.093} & \textbf{0.947} {\scriptsize $\pm$ 0.099} \\
        Animals 8 & 0.387 {\scriptsize $\pm$ 0.414} & 0.936 {\scriptsize $\pm$ 0.100} & \textbf{0.944} {\scriptsize $\pm$ 0.116} \\
        Fruits 8 & 0.027 {\scriptsize $\pm$ 0.122} & 0.768 {\scriptsize $\pm$ 0.396} & \textbf{0.977} {\scriptsize $\pm$ 0.044} \\
        \bottomrule
    \end{tabular}
    }
    \label{tab:entropy_gpt41}
\end{table}

\begin{table}[t]
    \centering
    \caption{Binary entropy (mean $\pm$ std) for \gptmini across prompt strategies. Best result per dataset in \textbf{bold}.}
    \resizebox{\textwidth}{!}{%
    \begin{tabular}{@{}lccc@{}}
        \toprule
        \textbf{Dataset} & \textbf{\basic} & \textbf{\explicit} & \textbf{\chainofthought} \\
        \midrule
        Mixed Categories & 0.950 {\scriptsize $\pm$ 0.224} & 0.998 {\scriptsize $\pm$ 0.010} & \textbf{1.000} {\scriptsize $\pm$ 0.000} \\
        \mcrae 8 & 0.600 {\scriptsize $\pm$ 0.457} & 0.975 {\scriptsize $\pm$ 0.058} & \textbf{0.995} {\scriptsize $\pm$ 0.014} \\
        GPT 8 & 0.729 {\scriptsize $\pm$ 0.408} & 0.986 {\scriptsize $\pm$ 0.043} & \textbf{1.000} {\scriptsize $\pm$ 0.000} \\
        \mcrae 16 & 0.619 {\scriptsize $\pm$ 0.462} & 0.937 {\scriptsize $\pm$ 0.110} & \textbf{0.984} {\scriptsize $\pm$ 0.068} \\
        \bigbench & 0.781 {\scriptsize $\pm$ 0.230} & 0.722 {\scriptsize $\pm$ 0.115} & \textbf{0.960} {\scriptsize $\pm$ 0.091} \\
        \things 8 & 0.872 {\scriptsize $\pm$ 0.206} & 0.876 {\scriptsize $\pm$ 0.255} & \textbf{0.963} {\scriptsize $\pm$ 0.093} \\
        Animals 8 & 0.826 {\scriptsize $\pm$ 0.246} & 0.855 {\scriptsize $\pm$ 0.169} & \textbf{0.956} {\scriptsize $\pm$ 0.066} \\
        Fruits 8 & 0.381 {\scriptsize $\pm$ 0.481} & \textbf{0.918} {\scriptsize $\pm$ 0.074} & 0.869 {\scriptsize $\pm$ 0.235} \\
        \bottomrule
    \end{tabular}
    }
    \label{tab:entropy_mini}
\end{table}

\para{The capability exists but is not default.}
The most striking result is the gap between \basic and \chainofthought prompting for \gptfull.
On Fruits~8, \basic achieves an entropy of only 0.027---the model generates questions like ``Is this a fruit?'' to which all items answer ``yes''---while \chainofthought achieves 0.977.
This 0.950 gap demonstrates that \gptfull possesses the knowledge to split homogeneous sets but does not deploy it without explicit reasoning instructions.
On structured sets (\mcrae 8, Mixed Categories), \chainofthought achieves perfect entropy of 1.000 with zero variance, indicating complete reliability.

\subsection{Perfect Split Rates}

\tabref{tab:perfect_splits} reports the percentage of trials achieving an exact $N/2$:$N/2$ split.

\begin{table}[t]
    \centering
    \caption{Perfect split rate (\%) for \gptfull. A perfect split means exactly $N/2$ items are answered ``yes'' and $N/2$ ``no.'' Best result per dataset in \textbf{bold}.}
    \begin{tabular}{@{}lccc@{}}
        \toprule
        \textbf{Dataset} & \textbf{\basic} & \textbf{\explicit} & \textbf{\chainofthought} \\
        \midrule
        Mixed Categories & 80.0 & \textbf{100.0} & \textbf{100.0} \\
        \mcrae 8 & 13.3 & 93.3 & \textbf{100.0} \\
        GPT 8 & 15.0 & \textbf{95.0} & 94.7 \\
        \mcrae 16 & 30.0 & 85.0 & \textbf{100.0} \\
        \bigbench & 0.0 & \textbf{86.7} & 13.3 \\
        \things 8 & 3.3 & 23.3 & \textbf{56.7} \\
        Animals 8 & 0.0 & 20.0 & \textbf{70.0} \\
        Fruits 8 & 0.0 & 25.0 & \textbf{65.0} \\
        \bottomrule
    \end{tabular}
    \label{tab:perfect_splits}
\end{table}

The perfect split rates sharpen the picture.
\chainofthought achieves 100\% perfect splits on Mixed Categories, \mcrae 8, and \mcrae 16---every single trial produces a question that divides the set exactly in half.
On harder sets, perfect splits drop to 56.7--70.0\%, indicating that while \chainofthought substantially improves performance, homogeneous sets remain challenging.
The exception is \bigbench, where \explicit (86.7\%) dramatically outperforms \chainofthought (13.3\%), as discussed in \secref{sec:setsize}.

\subsection{Statistical Analysis}

\para{Strategy comparison.}
Paired $t$-tests comparing prompt strategies for \gptfull show:
\begin{itemize}[leftmargin=*,itemsep=0pt,topsep=0pt]
    \item \basic vs.\ \chainofthought: $t = -16.27$, $p < 0.0001$, Cohen's $d = -1.76$ (very large effect)
    \item \basic vs.\ \explicit: $t = -15.01$, $p < 0.0001$, Cohen's $d = -1.61$ (very large effect)
    \item \chainofthought vs.\ \explicit: $t = 1.80$, $p = 0.073$, Cohen's $d = 0.19$ (not significant)
\end{itemize}

Both structured strategies massively outperform \basic prompting.
The difference between \chainofthought and \explicit is not statistically significant overall, though their relative performance varies by dataset.

\para{Model comparison.}
Comparing \gptfull and \gptmini under the \explicit strategy: $t = 2.24$, $p = 0.026$, Cohen's $d = 0.24$ (small effect favoring \gptfull).

\para{Dataset effects.}
One-way ANOVA across datasets under \explicit prompting for \gptfull yields $F = 6.35$, $p < 0.000001$, confirming significant differences in difficulty across set compositions.

\subsection{Set Composition Determines Difficulty}
\label{sec:difficulty}

Ranking datasets by \chainofthought perfect split rate for \gptfull reveals a clear difficulty gradient:

\begin{enumerate}[leftmargin=*,itemsep=0pt,topsep=0pt]
    \item \textbf{Mixed Categories} (100\%): Two distinct categories with an obvious boundary (\eg animals vs.\ vehicles).
    \item \textbf{\mcrae 8} (100\%): Two taxonomic categories (\eg birds and mammals).
    \item \textbf{GPT 8} (94.7\%): GPT-generated category pairs, similar to \mcrae.
    \item \textbf{\mcrae 16} (100\%): Larger structured sets remain manageable with reasoning.
    \item \textbf{\things 8} (56.7\%): Diverse random items without clear binary structure.
    \item \textbf{Animals 8} (70.0\%): All same category; must find sub-properties (\eg habitat, size).
    \item \textbf{Fruits 8} (65.0\%): All same category with limited discriminative features.
\end{enumerate}

The difficulty is driven by two factors.
First, \textbf{category homogeneity}: when all items share a category, no obvious binary split exists and the model must identify subtle properties (\eg ``name has more than six letters'').
Second, \textbf{feature availability}: fruits have fewer salient binary properties than animals (which can be split by habitat, diet, size, etc.), making Fruits the hardest 8-item dataset.

\subsection{The BigBench Anomaly: Set Size Effects}
\label{sec:setsize}

\bigbench (29 items) exhibits a striking reversal: \basic prompting achieves 0.992 entropy while \chainofthought achieves only 0.942, and \explicit achieves 86.7\% perfect splits versus \chainofthought's 13.3\%.

Two factors explain this anomaly.
First, with 29 diverse items spanning concrete objects, emotions, and abstract concepts, even a generic question like ``Is it a physical object?'' produces a near-even split by chance.
Second, chain-of-thought reasoning over 29 items is unreliable---the model loses count when enumerating properties and checking membership, leading to incorrect split predictions.
This suggests that \chainofthought is most effective for smaller sets (8--16 items) where the model can reliably track all items during reasoning.

\subsection{Smaller Models Have Better Default Behavior}

Under \basic prompting, \gptmini consistently outperforms \gptfull: 0.381 vs.\ 0.027 on Fruits, 0.826 vs.\ 0.387 on Animals, and 0.872 vs.\ 0.410 on \things.
\gptmini appears to interpret ``generate a yes/no question'' more naturally as a discriminative task, while \gptfull defaults to generating confirmation-style questions (\eg ``Is this a fruit?'' for a set of fruits).
This gap vanishes with explicit instructions---\gptfull catches up and slightly surpasses \gptmini under both \explicit and \chainofthought strategies.
